\documentclass[twocolumn,aps,pra,superscriptaddress,tightlines,nofootinbib]{revtex4-2}
\usepackage{hyperref}
\hypersetup{
	colorlinks=true,
	linkcolor=blue,         
	citecolor=blue,         
	filecolor=blue,      
	urlcolor=blue
}
\usepackage{url}
\usepackage{float}
\usepackage{placeins}
\usepackage{color}
\usepackage{graphicx,epsfig,colordvi,epstopdf}
\usepackage{bm}
\usepackage{mathtools}
\usepackage{mathdots}
\usepackage{physics}
\DeclareMathOperator*{\argmax}{arg\,max}
\DeclareMathOperator*{\argmin}{arg\,min}
\newcommand{\ra}[1]{\renewcommand{\arraystretch}{#1}}
\usepackage{amsmath}
\usepackage{amsfonts}
\usepackage{amssymb}

\newcommand\numberthis{\addtocounter{equation}{1}\tag{\theequation}}
\newcommand{\frank}[1]{\textcolor{blue}{#1}}

\begin{document}
	
	\setcounter{equation}{0}
	\setcounter{figure}{0}
	\setcounter{table}{0}
	\setcounter{page}{1}
	\makeatletter
	\renewcommand{\thesection}{S\arabic{section}}
	\renewcommand{\theequation}{S\arabic{equation}}
	\renewcommand{\thefigure}{S\arabic{figure}}
	
	\title{Supplemental Material: Mapping out phase diagrams with generative classifiers}
	\author{Julian Arnold}
	\affiliation{Department of Physics, University of Basel, Klingelbergstrasse 82, 4056 Basel, Switzerland}
	\author{Frank Sch\"afer}
	\affiliation{CSAIL, Massachusetts Institute of Technology, Cambridge, MA
		02139, USA}
	\author{Alan Edelman}
	\affiliation{CSAIL, Massachusetts Institute of Technology, Cambridge, MA
		02139, USA}
	\affiliation{Department of Mathematics, Massachusetts Institute of Technology, Cambridge, MA
		02139, USA}
	\author{Christoph Bruder}
	\affiliation{Department of Physics, University of Basel, Klingelbergstrasse 82, 4056 Basel, Switzerland}
	\date{\today}
	
	\maketitle
	
	\section{Additional details on phase-classification methods}
	\subsection{Assumptions when setting up phase-classification tasks}
	When setting up phase-classification tasks, we choose a uniform distribution over the set of parameters associated with each class, i.e., $P(\bm{\gamma}|y) = 1 / |\Gamma_{y}|$ for $\bm{\gamma}\in \Gamma_{y}$ and zero otherwise. This choice has been implicitly made via the definition of a loss function in previous works tackling the problem of mapping out phase diagrams with discriminative models (see Sec.~\ref{sec:discr_indicators}). Due to the application of Bayes' theorem within the generative approach, such probabilistic assumptions are brought to light. In principle, our framework allows for a free choice of $P(\bm{\gamma}|y)$. It represents our freedom to choose to what degree different points in parameter space are considered representative of the label $y$. Because the labeling procedure is different for the three schemes for mapping out phase diagrams discussed in the main text [Eqs. (3), (5), and (7) in the main text], its choice affects them differently. 
	
	In scheme 3 [Eq. (7) in the main text], each sampled value of the tuning parameter is considered its own class $\Gamma_{y} = \{ \bm{\gamma}_{y}\}$. As such, setting $P(\bm{\gamma}|y) = 1$ if $\bm{\gamma}  = \bm{\gamma}_{y}$ and zero otherwise is the only sensible choice. This corresponds to a special case of the uniform distribution.
	
	In scheme 2 [Eq. (5) in the main text], the sets $\Gamma_{1}^{(i)}(\bm{\gamma})$ and $\Gamma_{2}^{(i)}(\bm{\gamma})$ are each comprised of the $l$ points closest to $\bm{\gamma}$ in part 1 and 2 of the split parameter space. In the main text, we consider $l=1$ and each set is composed of a single point leaving only a single sensible choice as mentioned above for scheme 1. For $l > 1$, choosing $P(\bm{\gamma}|y)$ to be distinct from a uniform distribution would put different weights on different points within a given set. However, without additional system information, there is no compelling reason to do so. Moreover, note that for small $l$ and dense sampling of the parameter space, reweighting is expected to have a marginal effect since, in this case, $P(\bm{x}|\bm{\gamma})$ is similar for points $\bm{\gamma}$ within each of the two sets.
	
	In scheme 1 [Eq. (3) in the main text], the sets $\{ \Gamma_{y}\}_{y \in \mathcal{Y}}$ are chosen to be representative points within the $K$ distinct phases of the system. As such, scheme 1 already allows for flexibility in choosing points that are representative of a given label. While additional flexibility in the weighting of these points can be introduced, without additional system information there is no compelling reason to do so.
	
	\begin{figure}[bth!]
		\centering
		\includegraphics[width=0.9\linewidth]{fig_LBC_bias.pdf}
		\caption{Results for the cluster-Ising model [Eq.~(9) in the main text, $L=7$] at $h_{1}/J = 0.2$ measured using a Pauli-6 POVM. The set $\Gamma$ is composed of a uniform grid with 101 points. Indicator $I_{2}(\gamma)$ with $l = 101$ (i.e., considering global bipartitions of the one-dimensional parameter space) and in the presence (blue) or absence (black) of class imbalance. Estimated critical points (red dashed lines) are determined from maxima in $|\partial^2 \langle H \rangle_{\bm{\gamma}} /\partial \gamma_{2}^2|$. Here, the expected value involved in $I_{2}(\gamma)$ is computed exactly and the ground state is obtained via exact diagonalization.}
		\label{fig:LBC_imbalance}
	\end{figure}
	
	\subsection{Class imbalance}\label{sec_class_imb}
	In traditional discriminative machine learning (ML) tasks, class imbalance occurs when the number of samples representing each class within the data set differs. In phase-classification tasks, this can, for example, occur when one does not have access to an equal amount of data at each point in parameter space. Moreover, in scheme 1 (with indicator $I_{1}$) and 2 (with indicator $I_{2}$) discussed in the main text, an imbalance can also occur when sampling an \emph{equal} amount of data at each point in parameter space. In both cases, the imbalance arises because the sets $\{ \Gamma_{y}\}_{y \in \mathcal{Y}}$ are not of equal size. In scheme 1, this can be circumvented by choosing an equal number of points in parameter space to represent each phase. In scheme 2, data sets of unequal size occur as the bipartitions artificially divide the uniformly sampled one-dimensional parameter space. In previous works~\cite{van:2017,liu:2018,beach:2018,suchsland:2018,lee:2019,kharkov:2020,guo:2020,greplova:2020,bohrdt:2021,richter:2022,gavreev:2022,zvyagintseva:2022,guo:2023,schlomer:2023}, this class imbalance was not accounted for. As a result, the corresponding indicator $I_{2}$ will exhibit trivial local maxima ($I_{2}=1$) at the edges of the sampled region of the parameter space where the entire data is given the same label. If no phase transition is present, the indicator will thus exhibit a characteristic V-shape. In the presence of a phase transition, a W-shape is expected, where the location of the intermediate maximum corresponds to the predicted critical point. Such behavior is undesirable for an indicator of phase transitions, and can even lead to a failure to detect certain transitions (or misdetection). In Ref.~\cite{bohrdt:2021}, an attempt has been made to correct this pathological behavior. However, \citet{arnold:2022} have found that this procedure biases the transition point towards the center of the parameter range under consideration.
	
	In this work, we addressed class imbalances arising from sets $\{ \Gamma_{y}\}_{y \in \mathcal{Y}}$ of unequal size by choosing a uniform prior distribution $P(y)$. In particular, this removes the trivial local maxima at the edges of the sampled region of the parameter space previously encountered for the indicator $I_{2}$ and we instead expect a single peak in the presence of a phase transition. Figure~\ref{fig:LBC_imbalance} shows an example of transitions that are not detected using scheme 2 in the presence of class imbalance highlighting the importance of our proposed modification.
	
	\begin{figure}[bth!]
		\centering
		\includegraphics[width=0.99\linewidth]{fig_LBC_SI.png}
		\caption{Illustration for constructing local bipartitions for scheme 2 [Eq. (5) in the main text] of a two-dimensional parameter space ($d=2$) that is sampled (a) uniformly on a grid ($l=2$) or (b) non-uniformly ($l=3$). Crosses denote sampled points in parameter space $\bm{\gamma} \in \Gamma$. Sampled points belonging to one of the four sets $\Gamma_{1}^{(1)}(\bm{\gamma}^*)$,  $\Gamma_{2}^{(1)}(\bm{\gamma}^*)$, $\Gamma_{1}^{(2)}(\bm{\gamma}^*)$,  $\Gamma_{2}^{(2)}(\bm{\gamma}^*)$ are colored blue, green, pink, and orange, respectively. The corresponding splits of the parameter space are denoted by shaded colored regions.}
		\label{fig:LBC_bipartitions}
	\end{figure}
	
	
	\subsection{Scheme 2: How to choose local sets of points}
	In scheme 2 [Eq. (5) in the main text], at each sampled point $\bm{\gamma}^* = \left( \gamma_{1}^*,\gamma_{2}^*,\dots, \gamma_{d}^* \right) \in \Gamma$, the parameter space is split along each direction. For a given direction $1 \leq i \leq d$, this yields two sets, $\Gamma_{1}^{(i)}(\bm{\gamma}^*)$ and $\Gamma_{2}^{(i)}(\bm{\gamma}^*)$, each comprised of the $l$ points closest to $\bm{\gamma}^*$ in part 1 and 2 of the split parameter space, respectively. Using different norms for measuring this distance results in different local sets. Figure~\ref{fig:LBC_bipartitions} illustrates two distinct strategies. In this work, we choose the $l$ points $\bm{\gamma}$ closest to $\bm{\gamma}^*$ in each part of the split parameter space with respect to $|\gamma_{i} - \gamma_{i}^*|$, see Fig.~\ref{fig:LBC_bipartitions}(a). This is suitable for parameter spaces sampled on a uniform grid. In the case of non-uniformly sampled parameter spaces, one may instead consider a Euclidean distance $\norm{\bm{\gamma} - \bm{\gamma}^*}_{2}$, see Fig.~\ref{fig:LBC_bipartitions}(b). Note that in both cases, sets with a numbers of points lower than $l$ may arise at the edges of the parameter space. The resulting class imbalance is corrected for, see Sec.~\ref{sec_class_imb}.
	
	\subsection{Scheme 3: Dividing by the standard deviation}
	For the third phase-classification scheme discussed in the main text (with indicator $I_{3}$), we have proposed to divide the corresponding signal by its standard deviation. We find this modified indicator to be superior compared to its unmodified version
	\begin{align*}\label{eq:IPBM}
		I_{3}'(\bm{\gamma}) &= \sqrt{\sum_{i=1}^{d} \left(\frac{\partial \hat{\gamma}_{i}(\bm{\gamma})}{\partial \gamma_{i}}\right)^2},\\
		&=  \norm{ \mathbb{E}_{\bm{x} \sim P(\bm{x}|\bm{\gamma})}\left[\hat{\bm{\gamma}}(\bm{x}) \nabla_{\bm{\gamma}} \ln P(\bm{x}|\bm{\gamma}) \right]}_{2} \numberthis,
	\end{align*}
	which (up to a constant offset) has been considered in previous studies~\cite{schaefer:2019,greplova:2020,arnold:2021,arnold:2022}. As an example, Fig.~\ref{fig:fig_PBM_mod} compares the two indicators for the anisotropic Ising model. The indicator $I_{3}'$ shows four distinct peaks with the dominant two occurring within the two ordered phases (i.e., at larger $|J_{y}/k_{\rm B}T|$). Such spurious signals have also been observed in previous studies using indicators where the standard deviation is not taken into account~\cite{schaefer:2019,arnold:2022}. In contrast, the modified indicator $I_{3}$ we propose in this work shows two distinct peaks that qualitatively agree with the two critical points.
	
	\begin{figure}[tbh!]
		\centering
		\includegraphics[width=0.99\linewidth]{fig_PBM_mod.pdf}
		\caption{Results for the anisotropic Ising model on a square lattice ($L=20$) at $J_{x}/k_{\rm B}T = -0.325$. The indicator $I_{3}(\gamma)$ (black, left $y$-axis) and $I_{3}'(\gamma)$ (blue, right $y$-axis) are computed based on predictions $\hat{\bm{\gamma}}$ that are estimated element-wise from line scans across the entire two-dimensional phase diagram. The set $\Gamma$ is composed of a uniform grid with 60 points for each axis and $|\mathcal{D}_{\bm{\gamma}}| = 10^5 \; \forall \bm{\gamma} \in \Gamma$. Onsager's analytical solution for the phase boundaries is shown as a red dashed line.}
		\label{fig:fig_PBM_mod}
	\end{figure}
	
	\subsection{Discriminative cooperative networks}
	\indent In Ref.~\cite{liu:2018}, a scheme for mapping out phase diagrams in two-dimensional parameter spaces was proposed. Instead of a brute-force search of the entire parameter space for all possible phase boundaries, the predicted phase boundary is modeled as a parametrized curve using an active contour model called \emph{snake}~\cite{kass:1988}. This snake partitions the parameter space locally and is driven via internal forces that, e.g., prevent bending and stretching, as well as external forces aiming to minimize the overall classification error. In Ref.~\cite{liu:2018}, the external force is generated in an interplay between a guesser network and learner network, together called \emph{discriminative cooperative networks}, that are optimized via a joint cost function. The guesser provides labels for the data which the learner should reproduce. In each step, the guesser tries to provide a better set of labels based on the predictions of the learner to cooperatively minimize the classification error (i.e., the corresponding loss).
	
	\indent Let us first consider a one-dimensional parameter space. Following Ref.~\cite{liu:2018}, we use the following sigmoidal guesser
	\begin{equation}
		\tilde{P}_{\bm{\theta}_{G}}(1|\gamma) = \frac{1}{1+e^{(\gamma_{G} - \gamma)/\sigma_{G}}},
	\end{equation}
	where $\tilde{P}_{\bm{\theta}_{G}}(2|\gamma) = 1 - \tilde{P}_{\bm{\theta}_{G}}(1|\gamma)$. The guesser is characterized by two parameters $\bm{\theta}_{G}=(\gamma_{G},\sigma_{G})$, where $\gamma_{G}$ corresponds to the guessed transition point and $\sigma_{G}$ determines the sharpness of the transition. The learner is given by $\tilde{P}_{\bm{\theta}_{L}}(y|\bm{x})$. The guesser and learner are optimized jointly using a cross-entropy loss function
	\begin{widetext}
		\begin{equation}\label{eq:DCN}
			\mathcal{L}(\bm{\theta}) = -  \frac{1}{|\mathcal{Y}|}\sum_{y \in \mathcal{Y}} \frac{1}{|\Gamma_{y}|}\sum_{\gamma \in \Gamma_{y}}\frac{1}{|\mathcal{D}_{\gamma}|}\sum_{\bm{x} \in \mathcal{D}_{\gamma}} \tilde{P}_{\bm{\theta}_{G}}(y|\gamma) \log\left(\tilde{P}_{\bm{\theta}_{L}}(y|\bm{x})\right),
		\end{equation}
	\end{widetext}
	where $\bm{\theta} = \left( {\bm{\theta}_{G}}, {\bm{\theta}_{L}}\right)$. The parameters of the two networks can now be optimized using gradient descent on the loss function in Eq.~\eqref{eq:DCN}. This corresponds to a discriminative approach and has been used in Ref.~\cite{liu:2018}.
	
	\indent Moving to two-dimensional parameter spaces, one can utilize the snake as a parametrization for the guesser which is a discretized curve of linked nodes, $\bm{r}(s)=\left(\gamma_{1}(s),\gamma_{2}(s)\right)$, parametrized by $s \in [0,1]$ (assuming an open snake), see Fig.~\ref{fig:DCNs_LBC}(a). The snake moves in order to minimize its total energy $E_{\rm tot}=E_{\rm int} + E_{\rm ext}$. The internal energy
	\begin{equation}
		E_{\rm int} = \int_0^1\left(\alpha \norm{\frac{\partial \bm{r}}{\partial s}}_{2}^2+\beta\norm{\frac{\partial^2 \bm{r}}{\partial s^2}}_{2}^2\right) ds
	\end{equation}
	is introduced to make the snake smoother, where the hyperparameter $\alpha$ penalizes the stretching of the snake and $\beta$ penalizes its bending. The snake can sense its surroundings at each node through $2l$ sampled points perpendicular to the snake within a distance $l \sigma$ [see Fig.~\ref{fig:DCNs_LBC}(b)]. The overall guesser function is comprised of local guesser functions evaluated at each node that sense the direction perpendicular to the snake. The external energy $E_{\rm ext}$ of the snake corresponds to the overall loss obtained by summing the cross-entropy losses of all individual one-dimensional guessers [cf. Eq.~\eqref{eq:DCN}] and gives rise to an external force $-\delta E_{\rm ext}/\delta \mathbf{r}$ pointing perpendicular to the snake at each node.
	
	\indent In this work, we replace the learner network at each node by a corresponding generative classifier which makes updating the learner in each step obsolete. To construct the generative classifier, we derive the (empirically) optimal predictor based on the loss in Eq.~\eqref{eq:DCN}. Following the procedure outline in Sec.~\ref{sec:discr_indicators}, we obtain 
	\begin{equation}
		P_{\rm emp}(y|\bm{x}) = \frac{\sum_{\gamma \in \Gamma_{y}} \tilde{P}_{\bm{\theta}_{G}}(y|\gamma) \frac{M_{\gamma}(\bm{x})}{|\mathcal{D}_{\gamma}|} \frac{1}{|\mathcal{Y}|}\frac{1}{|\Gamma_{y}|}}{\sum_{y' \in \mathcal{Y}} \sum_{\gamma' \in \Gamma_{y'}} \tilde{P}_{\bm{\theta}_{G}}(y'|\gamma') \frac{M_{\gamma'}(\bm{x})}{|\mathcal{D}_{\bm{\gamma}'}|} \frac{1}{|\mathcal{Y}|}\frac{1}{|\Gamma_{y'}|}}.
	\end{equation}
	In the infinite-data limit, this converges to 
	\begin{equation}\label{eq:DCN_gen}
		P(y|\bm{x}) = \frac{\sum_{\gamma \in \Gamma_{y}} \tilde{P}_{\bm{\theta}_{G}}(y|\gamma)P(\bm{x}|\gamma)P(y)P(\gamma|y)}{\sum_{y'\in \mathcal{Y}} \sum_{\gamma' \in \Gamma_{y'}}\tilde{P}_{\bm{\theta}_{G}}(y'|\gamma')P(\bm{x}|\gamma')P(y)P(\gamma'|y')},
	\end{equation}
	with a uniform prior over the classes $P(y) = 1/|\mathcal{Y}|\; \forall y \in \mathcal{Y}$ and $P(\gamma|y) = 1/|\Gamma_{y}|$ if $\gamma \in \Gamma_{y}$ and zero otherwise. Thus, we can obtain a generative classifier by modeling $P(\bm{x}|\gamma)$ in Eq.~\eqref{eq:DCN_gen}. Note that the expression in Eq.~\eqref{eq:SI_optt} for binary labels is recovered as a special case.\\
	
	\begin{figure*}[htb!]
		\begin{center}
			\includegraphics[width=0.9\linewidth]{fig_SI_DCN.png}
			\caption{(a) Schematic illustration of the snake model. The blue circles represent the snake nodes (here 7). The normal (sensing) direction at each node is shown as a black-dashed line and the relevant $2l$ sampled points within a distance $2l \sigma$ are shown as black circles (here $l=2$). (b) Snake boundary for the anisotropic Ising model on a square lattice ($L=20$; $J_{x}/k_{\rm B}T,J_{y}/k_{\rm B}T\geq 0$) obtained using a generative classifier where the color bar denotes the training epoch. The snake consists of 20 nodes with $l=4$ and hyperparameters $\alpha = 0.002$, $\beta = 0.4$, $\kappa = 0.9$, and learning rates of $10^{-4}$ and $5 \times 10^{-4}$ associated with $E_{\rm int}$ and $E_{\rm ext}$, respectively. The set $\Gamma$ is composed of a uniform grid with 30 points for each axis and $|\mathcal{D}_{\bm{\gamma}}| = 10^5 \; \forall \bm{\gamma} \in \Gamma$. Onsager's analytical solution for the phase boundary is shown as a black-dashed line.
			}
			\label{fig:DCNs_LBC}
		\end{center}
	\end{figure*}
	
	In our implementation, we fix the width parameter $\sigma_{G}$ of the sigmoid guesser to be $\sigma_{G} = \sigma/10$, where $\sigma$ is the width parameter of the snake. We reduce $\sigma$ exponentially during the optimization according to 
	\begin{equation}
		\sigma_{k} = \sigma_{\rm end}+(\sigma_{\rm start}-\sigma_{\rm end}) \kappa^{k},
	\end{equation}
	where $k$ denotes the current epoch, $\sigma_{\rm start/end}$ are the start and end values, and $\kappa$ is the decay rate. We set $\sigma_{\rm end}=\Delta \gamma$, where $\Delta \gamma$ is the spacing between neighboring sampled points in parameter space and $\sigma_{\rm start} = 5 \sigma_{\rm end}$. This choice yields a strong gradient signal early on that becomes weaker but more accurate at later stages, which facilitates the convergence of the snake to the true underlying phase boundary. We construct models for $P(\bm{x}|\bm{\gamma})$ at points $\bm{\gamma}$ not contained within the initial sampled set $\Gamma$ from $\{  \tilde{P}(\bm{x}|\bm{\gamma})\}_{\bm{\gamma} \in \Gamma}$ using bilinear interpolation. Gradients of the energy with respect to the node positions are calculated using finite differencing. We minimize the snake's total energy using gradient-based optimization with Adam~\cite{kingma:2014}.\\ 
	
	The results obtained using this scheme with a generative learner for the anisotropic Ising model are shown in Fig.~\ref{fig:DCNs_LBC}(b). While the snake eventually finds the underlying phase boundary, there are several downsides to this scheme. First, the number of unique classification tasks that need to be solved over the course of the training is $N_{\rm epochs} \times N_{\rm nodes} = 400 \times 20 = 8000$. This is larger than the number of unique classification tasks that need to be solved within scheme 2 proposed in the main text which is given by $|\Gamma| \times 2 = 900 \times 2 = 1800$. Therefore, there is no gain in computation time compared to a brute-force search of the parameter space (as performed within scheme 2). The snake scheme also has a multitude of hyperparameters: $\alpha$, $\beta$, $\sigma_{\rm start}$, $\kappa$, $l$, learning rates, as well as the initial snake positioning and snake topology (e.g., whether the snake is open or closed and whether certain nodes should remain fixed throughout training). Thus, to run the snake scheme one has to perform hyperparameter tuning. This involves additional computational effort and/or prior knowledge of the underlying phase diagram. In particular, we find the results to depend heavily on $\alpha$, $\beta$, as well as the learning rates and initial position of the snake. Moreover, the scheme is not reliable in the presence of more than two phases as the snake will typically converge to one of the phase boundaries, missing the remaining ones [see Fig.~\ref{fig:DCNs_LBC}(b)]. To get around this, the scheme must be run multiple times with different snake initializations (possibly guided by prior knowledge of the phase diagram).\\  
	
	\section{Discriminative approach to mapping out phase diagrams}\label{sec:discr_indicators}
	The computation of all three indicators of phase transitions [Eqs. (3), (5), and (7) in the main text] boils down to solving classification tasks. In the main text, we discuss how such tasks can be solved in a generative manner. Here, we discuss how they can be solved in a discriminative manner. In this case, we look for the parameters $\bm{\theta}$ of a parametric model $\tilde{P}_{\bm{\theta}}(y|\bm{x})$ that minimize the following cross-entropy loss function
	\begin{equation}\label{eq:SI_x1}
		\mathcal{L}(\bm{\theta}) = - \frac{1}{|\mathcal{Y}|} \sum_{y \in \mathcal{Y}} \frac{1}{|\mathcal{D}_{y}|} \sum_{\bm{x} \in \mathcal{D}_{y}}  {\rm ln}\left(\tilde{P}_{\bm{\theta}}(y|\bm{x})\right),
	\end{equation}
	where $\mathcal{D}_{y} = \{ \bm{x} \in \mathcal{D}_{\bm{\gamma}} |\bm{\gamma} \in \Gamma_{y}\}$ is the relevant data set drawn from $P(\bm{x}|y)=1/|\Gamma_{y}| \sum_{\bm{\gamma} \in \Gamma_{y}} P(\bm{x}|\bm{\gamma})$ and $\mathcal{D}_{\bm{\gamma}}$ corresponding to a set of samples drawn from $P(\bm{x}|\bm{\gamma})$. Here, $|\mathcal{D}_{\bm{\gamma}}|$ is the same for all $\bm{\gamma} \in \Gamma_{y}$ (corresponding to the choice $P(\bm{\gamma}|y) = 1 / |\Gamma_{y}|$ for $\bm{\gamma}\in \Gamma_{y}$ and zero otherwise). Class imbalance is addressed by the rescaling factors $1/|\mathcal{D}_{y}|$ as is common in neural-network training~\cite{paszke:2019}. 
	
	Based on the trained parametric model, in scheme 1 we can estimate the central quantity of interest $P(y|\bm{\gamma})$ as 
	\begin{equation}\label{eq:SL_1}
		P(y|\bm{\gamma}) \approx \frac{1}{|\mathcal{D}_{\bm{\gamma}}|}\sum_{\bm{x} \in \mathcal{D}_{\bm{\gamma}}} \tilde{P}_{\bm{\theta}}(y|\bm{x})
	\end{equation}
	In scheme 2, the relevant quantity $p_{\rm err}$ can, for example, be estimated as
	\begin{equation}\label{eq:LBC_err_SI}
		p_{\rm err} \approx \frac{1}{2}\sum_{y \in \{1,2 \}} \frac{1}{|\mathcal{D}_{y}|}\sum_{\bm{x} \in \mathcal{D}_{y}} \left|y- \argmax_{y'} \tilde{P}_{\bm{\theta}}(y'|\bm{x})\right|.
	\end{equation}
	In scheme 3, the key quantity is $\hat{\bm{\gamma}}(\bm{x})$ which can be estimated as
	\begin{equation}\label{eq:PBM_pred_SI}
		\hat{\bm{\gamma}}(\bm{x}) \approx \sum_{y \in \mathcal{Y}} \tilde{P}_{\bm{\theta}}(y|\bm{x})\bm{\gamma}_{y}
	\end{equation}
	Alternatively, it may also be approximated directly using a parametric predictive model $\hat{\bm{\gamma}}_{\bm{\theta}}(\bm{x})$ that is trained to solve a regression task instead of a classification task (as has been done in previous works~\cite{schaefer:2019,greplova:2020,arnold:2021,arnold:2022}). In this case, the relevant loss function is
	\begin{equation}\label{eq:SI_x2}
		\mathcal{L}(\bm{\theta}) = \frac{1}{|\mathcal{Y}|} \sum_{y \in \mathcal{Y}} \frac{1}{|\mathcal{D}_{y}|} \sum_{\bm{x} \in \mathcal{D}_{y}}  \norm{\hat{\bm{\gamma}}_{\bm{\theta}}(\bm{x}) - \bm{\gamma}_{y}}_{2}^2.
	\end{equation}
	
	Above we have introduced the loss functions that are used to optimize parametric predictive models in the discriminative approach. In the following, we prove that in the infinite-data limit, an optimal predictive model, i.e., a model that minimizes the corresponding loss functions, yields Bayes optimal predictions. Analyzing its predictions yields the indicator signals discussed in the main text. Moreover, we show that an optimal discriminative classifier yields the same predictions as a generative classifier with $\tilde{P}(\bm{x}|\bm{\gamma})$ given by the empirical distribution obtained from the dataset $\mathcal{D}_{\bm{\gamma}}$.\\
	
	\indent For a given sample $\bm{x} \in \{x \in \mathcal{D}_{y}|y \in \mathcal{Y} \}$, we can determine the corresponding empirically optimal model prediction $P_{\rm emp}(y|\bm{x})$ by minimizing the loss function in Eq.~\eqref{eq:SI_x1} with respect to $\tilde{P}_{\bm{\theta}}(y|\bm{x})$ subjected to the equality constraint $\sum_{y\in\mathcal{Y}} \tilde{P}_{\bm{\theta}}(y|\bm{x}) - 1 = 0$. Using the method of Lagrange multipliers, the stationary points satisfy the following conditions
	\begin{equation}
		\lambda = \frac{1}{|\mathcal{Y}|} \frac{M_{y}(\bm{x})}{|\mathcal{D}_{y}|} \frac{1}{P_{\rm emp}(y|\bm{x})}\; \forall y \in \mathcal{Y},
	\end{equation}
	where $M_{y}(\bm{x})$ denotes the number of times $\bm{x}$ appears in $\mathcal{D}_{y}$ and $\lambda$ is the corresponding Lagrange multiplier. Together with the equality constraint, this yields
	\begin{equation}
		P_{\rm emp}(y|\bm{x}) = \frac{\frac{M_{y}(\bm{x})}{|\mathcal{D}_{y}|} \frac{1}{|\mathcal{Y}|}}{\sum_{y' \in \mathcal{Y}} \frac{M_{y'}(\bm{x})}{|\mathcal{D}_{y'}|} \frac{1}{|\mathcal{Y}|}}.
	\end{equation}
	Identifying the empirical distribution $P_{\rm emp}(\bm{x}|y) = M_{y}(\bm{x})/|\mathcal{D}_{y}|$ as well as the uniform prior over the classes $P(y) = 1/|\mathcal{Y}|\; \forall y \in \mathcal{Y}$, we have
	\begin{equation}\label{eq:SI_emp}
		P_{\rm emp}(y|\bm{x}) = \frac{P_{\rm emp}(\bm{x}|y)P(y)}{\sum_{y' \in \mathcal{Y}} P_{\rm emp}(\bm{x}|y')P(y')}.
	\end{equation}
	In the infinite-data limit $|\mathcal{D}_{y}| \rightarrow \infty$, we have $P_{\rm emp}(\bm{x}|y) \rightarrow P(\bm{x}|y)$. Thus, the predictions of the empirically optimal model converge to the predictions of a Bayes optimal model
	\begin{equation}\label{eq:SI_optt}
		P_{\rm emp}(y|\bm{x}) \rightarrow P(y|\bm{x}) = \frac{P(\bm{x}|y)P(y)}{\sum_{y' \in \mathcal{Y}} P(\bm{x}|y')P(y')},
	\end{equation}
	and we have recovered Eq.~(1) of the main text. Moreover, based on these optimal predictions, we obtain the optimal indicator of scheme 1 [Eq.~(3) of the main text]. To recover the optimal indicator of scheme 2, it remains to be shown that the estimated error rate in Eq.~\eqref{eq:LBC_err_SI} converges to the error rate as defined in Eq.~(4) of the main text in the case of an optimal discriminative model and infinite data. In this limit, Eq.~\eqref{eq:LBC_err_SI} transforms to 
	\begin{align*}\label{eq:perr_deriv}
		p_{\rm err} &= \frac{1}{2} \sum_{y \in \{1,2\}} \sum_{\bm{x}\in \mathcal{X}} P(\bm{x}|y)\left|y- \argmax_{y'} P(y'|\bm{x})\right|\\
		&= \frac{1}{2} \sum_{\substack{\bm{x}\in \mathcal{X}\\
				P(\bm{x}|1) \geq P(\bm{x}|2)}} P(\bm{x}|2) + \frac{1}{2} \sum_{\substack{\bm{x}\in \mathcal{X}\\
				P(\bm{x}|2) > P(\bm{x}|1)}} P(\bm{x}|1)\\
		&= \frac{1}{2} \left( 1- \frac{1}{2}\sum_{\bm{x}\in \mathcal{X}} \left|P(\bm{x}|1) - P(\bm{x}|2)\right|\right)\numberthis.
	\end{align*}
	For the last step, we use the fact that for any two probability distributions $p$ and $q$ over a discrete variable $x$, we have
	\begin{equation}
		\frac{1}{2}\sum_{x} \left|p(x)-q(x)\right| = 1-\sum_{\substack{x\\ q(x) > p(x)}} p(x) -\sum_{\substack{x\\ p(x) \geq q(x)}} q(x)\numberthis.
	\end{equation}
	Continuing with Eq.~\eqref{eq:perr_deriv}, we have
	\begin{align*}\label{eq:perr_deriv_2}
		p_{\rm err} &= \frac{1}{2} \left( \sum_{\bm{x}\in \mathcal{X}} P(\bm{x}) - \sum_{\bm{x}\in \mathcal{X}}P(\bm{x}) \left|P(1|\bm{x}) - P(2|\bm{x})\right|\right)\\
		&= \sum_{\bm{x}\in \mathcal{X}} P(\bm{x})\left( \frac{1}{2}\left( 1-  \left|P(1|\bm{x}) - P(2|\bm{x})\right| \right)\right)\\
		&= \sum_{\bm{x}\in \mathcal{X}}P(\bm{x}) \min\{P(1|\bm{x}), P(2|\bm{x})\},\numberthis
	\end{align*}
	where in step 2 we used Bayes' theorem $P(\bm{x}|y)/P(\bm{x})=P(y|\bm{x})/P(y)$ with $P(y)=1/2$. Plugging in the definition of
	\begin{align*}
		P(\bm{x}) &= \sum_{y \in \{ 1,2\}}P(\bm{x}|y)P(y) = \frac{1}{2}\sum_{y \in \{ 1,2\}} \sum_{\bm{\gamma} \in \Gamma }P(\bm{x}|\bm{\gamma})P(\bm{\gamma}|y)\\
		&= \frac{1}{2} \sum_{y \in \{ 1,2\}}\frac{1}{|\Gamma_{y}|} \sum_{\bm{\gamma} \in \Gamma_{y} }P(\bm{x}|\bm{\gamma})\numberthis
	\end{align*}
	in Eq.~\eqref{eq:perr_deriv_2} and defining $p_{\rm err}(\bm{x}) = \min\{P(1|\bm{x}), P(2|\bm{x})\}$ we recover the error rate in Eq.~(4) of the main text.\\
	
	\indent Following a similar minimization procedure for the loss function in Eq.~\eqref{eq:SI_x2}, we obtain
	\begin{align*}
		\hat{\bm{\gamma}}_{\rm emp}(\bm{x}) &= \sum_{y \in \mathcal{Y}} \frac{ \frac{M_{y}(\bm{x})}{|\mathcal{D}_{y}|} \frac{1}{|\mathcal{Y}|}}{\sum_{y' \in \mathcal{Y}} \frac{M_{y'}(\bm{x})}{|\mathcal{D}_{y'}|} \frac{1}{|\mathcal{Y}|}} \bm{\gamma}_{y}\\
		&= \sum_{y \in \mathcal{Y}} \frac{ P_{\rm emp}(\bm{x}|y) P(y)}{\sum_{y' \in \mathcal{Y}} P_{\rm emp}(\bm{x}|y') P(y')} \bm{\gamma}_{y},\numberthis
	\end{align*}
	which in the infinite-data limit converges to 
	\begin{equation}
		\hat{\bm{\gamma}}(\bm{x}) = \sum_{y\in \mathcal{Y}} \frac{P(\bm{x}'|y)P(y)}{\sum_{y'\in \mathcal{Y}} P(\bm{x}'|y')P(y')} \bm{\gamma}_{y} = \sum_{y\in \mathcal{Y}} P(y|\bm{x}) \bm{\gamma}_{y},
	\end{equation}
	corresponding to Eq.~(6) in the main text. Based on these predictions, we directly recover the indicator of scheme 3 [Eq.~(7) of the main text].\\
	
	\begin{figure*}[tbh!]
		\begin{center}
			\includegraphics[width=0.8\textwidth]{fig_SI_CI_timing.png}
			\caption{Root-mean-square error (RMSE) of estimated indicator of phase transitions $\tilde{I}$ as a function of the computation time associated with the discriminative (blue) and generative approach (black) to (a) scheme 1, (b) scheme 2, and (c) scheme 3 for the cluster-Ising model [Eq.~(9) in the main text, $L=7$] at $h_{1}/J = 0.2$. The total computation time is comprised of the time associated with data generation, training, and computation of the indicator based on the predictions, i.e., $\tilde{P}(y|\bm{x})$. For the data generation and construction of generative classifiers, we consider generative models based on matrix product states optimized via the density matrix renormalization group algorithm. The indicators $\tilde{I}$ are estimated from datasets $\{ \mathcal{D}_{\bm{\gamma}} \}_{\bm{\gamma} \in \Gamma}$ of various size. The reference indicators $I$ are computed exactly based on ground states obtained via exact diagonalization. The set $\Gamma$ is composed of a uniform grid with 101 points. In scheme 1, $\Gamma_{1} = \{ -1.5 \; h_{2}/J \}$, $\Gamma_{2} = \{ -0.03 \; h_{2}/J \}$, and $\Gamma_{3} = \{ 1.47 \; h_{2}/J \}$. In scheme 2, we consider $l=2$. The discriminative approach to scheme 3 is performed according to Eq.~\eqref{eq:SI_x2}. For the generative approach, 10 independent runs with distinct datasets are shown. For the discriminative approach, we considered feedforward NNs of various sizes (ranging from a single hidden layer with 64 nodes to five hidden layers with 64 nodes each; containing rectified linear units as activation functions), different dataset sizes $|\mathcal{D}_{\bm{\gamma}}|$ (ranging from 10 to $2 \times 10^{5}$), learning rates (ranging from $5 \times 10^{-5}$ to $5 \times 10^{-1}$), and number of training epochs (ranging from 1 to $5 \times 10^{4}$). Each blue curve corresponds to a training run with a fixed choice of hyperparameters and we keep track of the best RMSE throughout training. The NNs are implemented using Flux~\cite{innes:2018} in \texttt{Julia}, where the weights and biases are optimized via gradient descent with Adam~\cite{kingma:2014}. Gradients are calculated using backpropagation~\cite{baydin:2018}. The NN inputs are composed of six integers per qubit encoding the POVM element as well as the corresponding measurement outcome and are standardized before training.}
			\label{fig:CI_timing}
		\end{center}
	\end{figure*}
	
	\indent Note that using
	\begin{align*}
		P_{\rm emp}(\bm{x}|y) &= \frac{1}{|\Gamma_{y}|} \sum_{\bm{\gamma} \in \Gamma_{y}} \frac{M_{\bm{\gamma}}(\bm{x})}{|\mathcal{D}_{\bm{\gamma}}|}\\
		&= \frac{1}{|\Gamma_{y}|} \sum_{\bm{\gamma} \in \Gamma_{y}} P_{\rm emp}(\bm{x}|\bm{\gamma})\numberthis,
	\end{align*}
	Eq.~\eqref{eq:SI_emp} coincides with Eq. (2) in the main text with $P(\bm{x}|\bm{\gamma})$ replaced by $P_{\rm emp}(\bm{x}|\bm{\gamma})$. That is, the predictions of an optimal discriminative model are equivalent to the ones of a generative model with $\tilde{P}(\bm{x}|\bm{\gamma}) = P_{\rm emp}(\bm{x}|\bm{\gamma})$.
	
	\section{Comparison between discriminative and generative approach}
	In this section, we compare the computational cost associated with the discriminative and generative approaches to mapping out phase diagrams. Let $t^{\rm gen}_{\rm eval}$ and $t^{\rm discr}_{\rm eval}$ be the times corresponding to evaluating the generative model $\tilde{P}(\bm{x}|\bm{\gamma})$ or discriminative model $\tilde{P}(y|\bm{x})$ for a given $\bm{x}$, respectively. These times depend heavily on the choice of model, implementation, and hardware. Here, the reported computation times were assessed in \texttt{Julia} (version 1.8.2) on a single 3.70 GHz Intel Xeon W-2135 CPU. For the non-parametric generative model of the anisotropic Ising model ($L=20$), the evaluation time is negligible $t^{\rm gen}_{\rm eval} \approx 40$ ns. In the case of the matrix product state-based generative models for the ground states of the cluster-Ising Hamiltonian, the evaluation time $t^{\rm gen}_{\rm eval}$ ranges from $\approx 96$~${\rm \mu s}$ for $L=7$ (on average across the sampled parameter space) to $\approx 3$~ms for $L=71$. For comparison, the evaluation time associated with a feedforward NN is $t^{\rm discr}_{\rm eval} \approx 7$~${\rm \mu s}$ in the case of a single hidden layer containing a single node and, for example, $t^{\rm discr}_{\rm eval} \approx 2$~ms in the case of five hidden layers containing 1024 nodes each (assuming only a single input and output node for simplicity).
	
	\indent The computation time of a single prediction via the generative approach, i.e., computation of $\tilde{P}_{\rm gen}(y|\bm{x}) \; \forall y \in \mathcal{Y}$ for a given sample $\bm{x}$ via Eq. (2) in the main text, scales as $t^{\rm gen}_{\rm pred} =  t^{\rm gen}_{\rm eval} | \Gamma_{\mathcal{Y}}| $, where $\Gamma_{\mathcal{Y}}= \{ \bm{\gamma} \in \Gamma_{y}| y \in \mathcal{Y}\}$. In the main text, in scheme 1 we have $| \Gamma_{\mathcal{Y}}| = K = 5$ and $3$, in scheme 2, $| \Gamma_{\mathcal{Y}}| = l = 1$, and in scheme 3, $| \Gamma_{\mathcal{Y}}| = 60$ and $101$, in the case of the anisotropic Ising and cluster-Ising model, respectively. The time associated with computing $\tilde{P}_{\rm discr}(y|\bm{x})$ via the discriminative approach is $t^{\rm discr}_{\rm pred} = t^{\rm discr}_{\rm eval}$. Given a model $\tilde{P}(y|\bm{x})$ (either discriminative or generative), the time associated with computing an indicator $I(\bm{\gamma}) \; \forall \bm{\gamma} \in \Gamma$ scales as $t_{\rm pred} |\Gamma| |\mathcal{D}_{\bm{\gamma}}|$ in case of scheme 1 and 3, and $t_{\rm pred} |\Gamma| |\mathcal{D}_{\bm{\gamma}}||\Gamma_{y}|$ in case of scheme 2. Based on this, for the anisotropic Ising model, the generative approach yields a speedup in the indicator computation of \textit{at least} (assuming a feedforward NN of minimal size as the discriminative model) a factor of $ t_{\rm pred}^{\rm discr}/t_{\rm pred}^{\rm gen} \approx 3$, $35$, and $87$ in the case of scheme 3, 1, and 2, respectively.
	
	\indent In scenarios where the generative model acts as a data source, such as in the two cases considered in the main text, the discriminative approach is guaranteed to have an overhead with respect to the generative approach due to training. We can estimate this overhead, i.e., the computation time required until the first evaluation of $\tilde{P}_{\rm discr}(y|\bm{x})$, to scale as $t^{\rm discr}_{\rm eval} | \Gamma_{\mathcal{Y}}| |\mathcal{D}_{\bm{\gamma}}|N_{\rm epochs}$. Note that to predict the indicator in scheme 2, $d |\Gamma|$ such discriminative models need to be trained (one for each bipartition at each point in parameter space), making the overhead scale accordingly. The number of training epochs is denoted by $N_{\rm epochs}$. In practice, $N_{\rm epochs}$ may range from $\mathcal{O}(10^2) - \mathcal{O}(10^4)$ and above. In cases where $\tilde{P}(\bm{x}|\bm{\gamma}) \approx P(\bm{x}|\bm{\gamma})$, for the discriminative approach to yield a predictive model of comparable quality, one expects it to be highly expressive (resulting in a high $t^{\rm discr}_{\rm eval}$), trained for a long time, and with a large dataset, resulting in a large overhead overall. The requirement of a larger dataset comes with an additional cost (with respect to the generative approach) associated with sample generation. 
	
	Because of this overhead, the generative approach can be more computationally efficient than the discriminative approach \textit{even if} $t_{\rm pred}^{\rm gen} > t_{\rm pred}^{\rm discr}$. This is illustrated in Fig.~\ref{fig:CI_timing} for the case of the cluster-Ising model. For scheme 2 [Fig.~\ref{fig:CI_timing}(b)], the overhead of the discriminative approach is so large that the generative approach is more efficient irrespective of the error tolerance. In particular, in scheme 2, for small $l$ it is expected to be difficult to achieve accurate results using the discriminative approach with a small dataset, because of the small differences in the underlying distributions. For schemes 1 and 3 [Fig.~\ref{fig:CI_timing}(a) and (c)], there exists a crossover point in terms of the error of the estimated indicator below which the generative approach is computationally more efficient. Note that hyperparameter tuning is required for the discriminative approach to yield accurate estimates of the indicator, which is not accounted for in the overall computation time. In contrast, using the generative approach, the estimated indicator can be systematically improved by sampling more data.
	
	\section{Mapping out phase diagrams of classical equilibrium systems}
	\subsection{Optimal lossless compression of state space}
	In this section, we prove that to map out phase diagrams of a large class of classical equilibrium systems by computing one of the three indicators discussed in the main text, it suffices to model $\{ P(\bm{X}|\bm{\gamma})\}_{\bm{\gamma} \in \Gamma}$, where $\bm{X}$ is the corresponding sufficient statistic.
	
	Consider probability distributions of the exponential family
	\begin{equation}\label{eq:SM_exp_fam}
		P_{\rm exp}(\bm{x}|\bm{\gamma}) = h(\bm{x})\exp(\sum_{i=1}^{d} \eta_{i}(\bm{\gamma}) X_{i}(\bm{x})-A(\bm{\gamma})),
	\end{equation}
	where $h(\bm{x})\geq 0$ is the carrier measure, $\bm{X}(\bm{x})$ is a sufficient statistic, $\bm{\eta}(\bm{\gamma})$ are the natural parameters, and $A(\bm{\gamma})$ is the log-partition function. The statistic $\bm{X}(\bm{x})$ is a \emph{minimal} sufficient statistic for $\bm{\gamma}$ if the set of allowed natural parameters $\bm{\eta}(\bm{\gamma})$ spans a $d$-dimensional space~\cite{casella:2002}. The distribution over the sufficient statistic is given by
	\begin{align*}
		P_{\rm exp}(\bm{X}'|\bm{\gamma}) &= \sum_{\bm{x} \in \mathcal{X}} P_{\rm exp}(\bm{x}|\bm{\gamma}) \delta(\bm{X}(\bm{x}) - \bm{X}')\\
		&= g(\bm{X}')\exp(\bm{\eta}(\bm{\gamma}) \cdot \bm{X}' - A(\bm{\gamma}))\numberthis,
	\end{align*}
	where $g(\bm{X}') = \sum_{\bm{x} \in \mathcal{X}} h(\bm{x})\delta(\bm{X}(\bm{x}) - \bm{X}')$. We have 
	\begin{equation}\label{eq:SM_der}
		P_{\rm exp}(\bm{x}|\bm{\gamma})/P_{\rm exp}(\bm{X}(\bm{x})|\bm{\gamma})=h(\bm{x})/g(\bm{X}(\bm{x})).
	\end{equation}
	\begin{widetext}
		Crucially, this ratio is independent of $\bm{\gamma}$. Using Eq.~\eqref{eq:SM_der}, the conditional probability of a label in phase-classification tasks [Eq.~(2) in the main text] can be written as
		\begin{equation}\label{eq:Bayes_2_SM}
			P(y|\bm{x}) =  \frac{\frac{1}{|\Gamma_{y}|}\sum_{\bm{\gamma} \in \Gamma_{y}} P_{\rm exp}(\bm{x}|\bm{\gamma})}{\sum_{y' \in \mathcal{Y}}\frac{1}{|\Gamma_{y'}|}\sum_{\bm{\gamma}' \in \Gamma_{y'}} P_{\rm exp}(\bm{x}|\bm{\gamma}')}
			=  \frac{\frac{1}{|\Gamma_{y}|}\sum_{\bm{\gamma} \in \Gamma_{y}} P_{\rm exp}(\bm{X}(\bm{x})|\bm{\gamma})}{\sum_{y' \in \mathcal{Y}}\frac{1}{|\Gamma_{y'}|}\sum_{\bm{\gamma}' \in \Gamma_{y'}} P_{\rm exp}(\bm{X}(\bm{x})|\bm{\gamma}')}
			= P(y|\bm{X}(\bm{x})).
		\end{equation}
		Thus,
		\begin{align*}\label{eq:SM_der_1}
			P(y|\bm{\gamma}) &= \mathbb{E}_{\bm{x} \sim P_{\rm exp}(\bm{x}|\bm{\gamma})}\left[ P(y|\bm{x}) \right] =  \sum_{\bm{x} \in \mathcal{X}} P_{\rm exp}(\bm{x}|\bm{\gamma})  P(y|\bm{x}) = \sum_{\bm{x} \in \mathcal{X}}  P_{\rm exp}(\bm{X}(\bm{x})|\bm{\gamma})  P(y|\bm{X}(\bm{x})) h(\bm{x})/g(\bm{X}(\bm{x}))\\
			&= \sum_{\bm{X}' \in \mathcal{X}_{\rm suff}}\sum_{\bm{x} \in \mathcal{X}} P_{\rm exp}(\bm{X}'|\bm{\gamma})  P(y|\bm{X}') \delta(\bm{X}(\bm{x})-\bm{X}') h(\bm{x})/g(\bm{X}')= \sum_{\bm{X} \in \mathcal{X}_{\rm suff}} P_{\rm exp}(\bm{X}|\bm{\gamma})  P(y|\bm{X}) = \mathbb{E}_{\bm{X} \sim P_{\rm exp}(\bm{X}|\bm{\gamma})}\left[  P(y|\bm{X})\right] \numberthis, \\
		\end{align*}
	\end{widetext}
	where $\mathcal{X}_{\rm suff}$ is the state space (without duplicates) associated with the sufficient statistic $\mathcal{X}_{\rm suff} = \{\bm{X}(\bm{x}) | \bm{x} \in \mathcal{X} \}$. Similarly, using Eq.~\eqref{eq:Bayes_2_SM}, we have
	\begin{align*}
		\hat{\bm{\gamma}}(\bm{x}) &=  \sum_{y \in \mathcal{Y}} P(y|\bm{x})\bm{\gamma}_{y} = \sum_{y \in \mathcal{Y}} P(y|\bm{x})\bm{\gamma}_{y} P(y|\bm{X}(\bm{x}))\bm{\gamma}_{y}\\
		&= \hat{\bm{\gamma}}\left(\bm{X}(\bm{x})\right) \numberthis,
	\end{align*}
	and
	\begin{align*}\label{eq:SM_der_2}
		\hat{\bm{\gamma}}(\bm{\gamma}) &= \mathbb{E}_{\bm{x} \sim P_{\rm exp}(\bm{x}|\bm{\gamma})}\left[ \hat{\bm{\gamma}}(\bm{x})\right] = \sum_{\bm{x} \in \mathcal{X}} P_{\rm exp}(\bm{x}|\bm{\gamma})\hat{\bm{\gamma}}(\bm{x})\\
		&=  \sum_{\bm{X} \in \mathcal{X}_{\rm suff}} P_{\rm exp}(\bm{X}|\bm{\gamma}) \hat{\bm{\gamma}}(\bm{X}) = \mathbb{E}_{\bm{X} \sim P_{\rm exp}(\bm{X}|\bm{\gamma})}\left[  \hat{\bm{\gamma}}(\bm{X})\right]\numberthis.
	\end{align*}
	The relevant indicators [Eqs. (6), (8), and (10)] of the three phase-classification methods presented in the main text can be straightforwardly computed based on the quantities in Eqs.~\eqref{eq:SM_der_1}, and~\eqref{eq:SM_der_2}, which are expressed solely in terms of the distribution over the sufficient statistic $P(\bm{X}|\bm{\gamma})$ (instead of the full distribution over $\bm{x}$).
	
	In the main text, we considered classical systems with dimensionless Hamiltonians of the form
	\begin{equation}\label{eq:exp}
		\mathcal{H}(\bm{x},\bm{\gamma}) = \sum_{i=1}^{d} \gamma_{i} X_{i}(\bm{x}),
	\end{equation}
	at equilibrium with a large thermal reservoir. In this case, the probability of finding the system in state $\bm{x} \in \mathcal{X}$ is given by
	\begin{equation}\label{eq:Boltzmann_SM}
		P(\bm{x}| \bm{\gamma}) = e^{-\mathcal{H}(\bm{x},\bm{\gamma})}/Z(\bm{\gamma}).
	\end{equation}
	Such Boltzmann distributions belong to the exponential family [Eq.~\eqref{eq:SM_exp_fam}] with $h(\bm{x}) = 1$, $A(\bm{\gamma}) = {\rm ln}\,Z(\bm{\gamma})$, and $\eta_{i}(\bm{\gamma}) = \gamma_{i}$. Thus, as a special case of the above, it follows that phase diagrams of such systems can be mapped out given distributions over the sufficient statistic $\{P(\bm{X}|\bm{\gamma})\}_{\bm{\gamma} \in \Gamma}$ (instead of the full distributions over $\bm{x}$).
	
	\subsection{Utilizing knowledge of symmetries}
	The sufficient statistic allows for an optimal lossless compression of the state space $\mathcal{X}$. However, it requires knowledge of the underlying Hamiltonian. In the following, we demonstrate how one can achieve a lossless compression of the state space without explicitly knowing the underlying Hamiltonian, but only utilizing knowledge of its symmetries. Let us consider a symmetry operation $S: \mathcal{X} \rightarrow \mathcal{X}$ that leaves the energy of the system invariant $\mathcal{H}\left(S(\bm{x})\right) = \mathcal{H}(\bm{x}) \;\forall \bm{x} \in \mathcal{X}$. Considering dimensionless Hamiltonians of the form given in Eq.~\eqref{eq:exp}, we have
	\begin{equation}
		\sum_{i=1}^{d} \gamma_{i} \left[X_{i}(\bm{x}) - X_{i}\left(S(\bm{x})\right)\right] = 0
	\end{equation}
	for all allowed parameters $\bm{\gamma}$. Assuming that this set spans a $d$-dimensional space, we have $\bm{X}(\bm{x}) = \bm{X}\left(S(\bm{x})\right)\; \forall \bm{x} \in \mathcal{X}$. That is, the sufficient statistic is also invariant under any symmetry operation $S$. Thus, we can perform a lossless compression by adopting a representation that is unique for all samples related by symmetry operations. Let us denote the associated state space by $\mathcal{X}_{\rm symm}$. Note that while this compression is lossless, it is not necessarily optimal, i.e., $|\mathcal{X}_{\rm symm}| \geq |\mathcal{X}_{\rm suff}| $.
	
	\section{Data generation}
	Here, we provide further details on the data-generation process for the (classical) anisotropic Ising model and the (quantum) cluster-Ising model.
	\subsection{Anisotropic Ising model}
	Given a set of parameters $\bm{\gamma}=(J_{x}/k_{\rm B}T,J_{y}/k_{\rm B}T)$, we use the Metropolis-Hastings algorithm to sample spin configurations from the corresponding Boltzmann distribution. The lattice is updated by drawing a random spin, which is flipped with probability ${\rm min}(1, e^{-\Delta E/k_{\rm B}T})$, where $\Delta E$ is the energy difference resulting from the considered flip. After a thermalization period of $10^5$ lattice sweeps, we collect $10^{5}$ samples. We treat each quadrant of the phase diagram separately. In each quadrant, we initialize the system in one of the two corresponding ground states and increase $\gamma_{1}$ at constant $\gamma_{2}$. When increasing $\gamma_{1}$, we use the spin configuration from the preceding value as an initial condition for the Markov chain. When increasing $\gamma_{2}$, we reset the system to the corresponding ground state.
	
	\subsection{Cluster-Ising model}
	Given a set of parameters $\bm{\gamma}=(h_{1}/J,h_{2}/J)$, we compute ground states of the corresponding cluster-Ising Hamiltonian [Eq.~(9) in main text] with $L = 71$ spins using the finite-size density matrix renormalization group (DMRG) with a maximum bond dimension of 150. We utilize the ITensor package~\cite{itensor} in \texttt{Julia}.
	
	\bibliography{refs.bib}
\end{document}